$k(y)$-向量空间 $(\Omega^1_{Y/k})_y$ 的维数为 $n+1$. 将标量从 $k(y)$ 扩张到其代数闭包, 并在 $y$ 处完备化后检验即可; 此时结论显然. 因而 $Y$ 中 $y$ 的一个邻域是某个维数为 $n+1$ 的光滑 $k$-概形 $Z$ 的闭子概形. 此外, 将 $Y$ 和 $Z$ 换成适当的 $y$ 的邻域后, $\mathcal O_Z$ 中定义 $Y$ 的理想为主理想(论证同上), 设其方程为 $f$.

\subsubsection*{引理 1.2.7.}
\underline{假设 $y$ 是 $Y$ 的非退化二次点, 且 $k(y)/k$ 为纯不可分扩张. 则在 $y$ 的一个邻域中, $J^n_{Y/k}$ (1.1.1)由一个元素生成, 支集为 $y$, 在 $k$ 上秩为一. 特别地, $k(y)=k$.}

转到 $k$ 的代数闭包并在 $y$ 处完备化后证明第一个断言即可. 此时直接计算即可得出: 采用 1.2.1 的记号, Nakayama 引理表明 $\partial f/\partial x_i$(从而更有 $f$ 及 $\partial f/\partial x_i$)生成极大理想.

现在对非退化的 $y$ 证明 1.2.6. 可以而且将假设 $k(y)=k$, $Y$ 在光滑概形 $Z$ 中由一个方程 $f$ 定义, 且有将 $y$ 映到 $0$ 的平展态射
\[
x:Z\longrightarrow\mathbb A_k^{n+1}
\]
(一个“局部坐标系 $x_i$”). 于是存在 $n+1$ 元非退化二次型 $Q$, 使在 $y$ 处有
\[
f=Q(x_i)+\text{次数为 }>2.
\]
设 $Y_0$ 如 1.2.3 中那样由 $Q$ 定义.\footnote{印本此处引用 4.3; 刚使用的正规形实际为例 1.2.3.}

\underline{在 $Y$ 上}, 系统 $(x_i)$ 为方程 $Q(x_i)=0$ 的一个“近似解”: 若 $\mathfrak m$ 是定义 $y$ 的理想, 则
\[
Q(x_i)\in\mathfrak m^3\qquad\text{在 }Y.
\]
由于 $Q$ 非退化, $\mathbb A_k^{n+1}$ 上由 $\partial Q/\partial x_i$ 生成的理想就是 $0$ 处的极大理想, 它在 $x=(x_i):Y\to\mathbb A_k^{n+1}$ 下的逆像为 $\mathfrak m$. 由隐函数定理 1.1.2(取 $a=\mathfrak m$ 及 $\delta'=J$), Hensel 化 $Y_{(y)}$ 上存在函数 $x'_i$, 模 $\mathfrak m^2$ 与 $x_i$ 同余, 且满足 $Q(x'_i)=0$. 它们给出所求的 $Y_{(y)}$ 与 $Y_{0(y_0)}$ 之间的同构.

\subsubsection*{引理 1.2.8.}
\underline{假设 $y$ 为退化的普通二次点($p=2$ 且 $n$ 为偶数), 并假设 $k(y)/k$ 为纯不可分扩张. 则在 $y$ 的一个邻域中, $T^1_{Y/k}$ 由一个元素生成, 支集为 $y$, 在 $k$ 上秩为二. 特别地, 或者 $k(y)=k$, 或者对某个 $a\in k$, $a\notin k^2$, $k(y)$ 同构于 $k(\sqrt a)$.}

采用 1.2.7 的方法. 用 1.2.1 的记号, 只需检验
\[
\dim_k\left(
k[[x_1,\ldots,x_{n+1}]]\big/
  (f,\partial f/\partial x_i)
\right)=2.
\]
在适当坐标 $x_0,x_1,\ldots,x_{2m}$ 下, 形式幂级数 $f$ 可写为
\[
f=x_0^2+\sum_{i=1}^m x_ix_{i+m}+R,
\qquad\operatorname{ord}(R)\geq3.
\]
设 $\mathfrak n'$ 为由 $f$ 和 $\partial f/\partial x_i$ 生成的理想, $\mathfrak n$ 为由 $x_0^2$ 及 $i\ne0$ 的 $x_i$ 生成的理想, $\mathfrak q$ 为极大理想. 我们证明 $\mathfrak n=\mathfrak n'$. 显然 $\mathfrak n'\subset\mathfrak n$, 由 Nakayama 引理, 只需检验
\[
\mathfrak n'+\mathfrak n\mathfrak q=\mathfrak n.
\]
模 $\mathfrak n\mathfrak q$, 有 $f\equiv x_0^2$; 模 $\mathfrak n\mathfrak q+(x_0^2)$, 有
\[
\frac{\partial f}{\partial x_i}\equiv x_{i+m},
\qquad
\frac{\partial f}{\partial x_{i+m}}\equiv x_i
\qquad(1\leq i\leq m),
\]
这就证明了断言.

在 1.2.8 的假设下, 若 $k=k(y)$, 则以 $a$ 表示 $k$ 中任意一个平方.

\subsubsection*{引理 1.2.9.}
\underline{设 $V$ 为概形 $S$ 上的向量丛, $T$ 为 $\mathbb A_S(V)$ 中在 $S$ 上局部自由秩为二的闭子概形. 则 $T$ 包含在 $\mathbb A_S(V)$ 中唯一的一条直线内.}

此问题对于 $S$ 上的 fppf 拓扑是局部的, 因此可假设 $T/S$ 有截面. 经过平移, 可假设 $T$ 包含零截面. 设 $\mathfrak n$ 为 $T$ 中定义 $0$ 的理想; 它在 $S$ 上局部自由秩为一. $\mathbb A_S(V)$ 上每个齐次线性形式 $v\in V$ 都定义 $\mathfrak n$ 的一个元素, 从而得到满射 $V\twoheadrightarrow\mathfrak n$. 其核定义所求直线.

\subsubsection*{1.2.10.}
特别地, 若 $A$ 是特征 $2$ 的域 $k$ 上仿射空间 $\mathbb A_k^{n+1}$ 的纯不可分秩二闭子概形, 则存在 $a\in k$, 使在适当的仿射坐标下, $A$ 的方程为
\[
x_0^2-a=0,
\qquad x_i=0\quad(1\leq i\leq n).
\]
只需取在 $A$ 所确定的直线上消失的 $x_i$.

\subsubsection*{1.2.11.}
现在对退化的 $y$($p=2$ 且 $n=2m$)证明 1.2.6. 可以而且将假设 $k(y)=k(\sqrt a)$, $Y$ 在光滑概形 $Z$ 中由一个方程 $f$ 定义, 且在 $y$ 处存在一个局部坐标系
\[
x=(x_i)_{0\leq i\leq2m}:
Z\longrightarrow\mathbb A_k^{2m+1}
\]
将 $y$ 映到闭点 $(\sqrt a,0,\ldots,0)$, 使在 $y$ 的一个邻域中有
\[
J^n_{Y/k}\simeq
\mathcal O_Y/(x_0^2-a,x_1,\ldots,x_{2m}).
\]
将 $f$ 换成适当的倍数后, 有
\[
f=(x_0^2-a)+
 \sum_{1\leq i\leq j\leq2m}a_{ij}x_ix_j+R,
\]
其中
\begin{enumerate}[label=(\alph*),leftmargin=2.6em]
\item $2m$ 元二次型 $\sum a_{ij}x_ix_j$ 非退化;
\item 将标量扩张到 $\bar k$ 后, $R$ 关于 $x_0-\sqrt a$ 和 $x_i$ 的阶至少为 $3$.
\end{enumerate}
设 $Y_0$ 如 1.2.4 中那样由下式定义:
\[
Q=x_0^2-a+\sum a_{ij}x_ix_j.
\]

\underline{在 $Y$ 上}, 系统 $(x_i)$ 仍可视作 $Q(X_i)=0$ 的近似解. 这一次, 由 $\partial Q/\partial X_i$ 生成的理想在 $x$ 下的逆像 $\delta$ 满足 $\mathcal O_Y/\delta\simeq J^n_{Y/k}$; 在 $y$ 处它为
\[
\delta=(x_0^2-a,(x_i)_{1\leq i}).
\]
设 $\mathfrak q$ 为定义 $y$ 的极大理想. 为了像非退化情形那样应用隐函数定理得出结论, 尚须检验
\begin{equation}\tag{1.2.11.1}
Q(x_i)\equiv0\pmod{\delta^2\mathfrak q}
\qquad\underline{\text{在 }Y}.
\end{equation}
设 $\bar{\mathfrak q}$ 为将标量扩张到 $k$ 的代数闭包后定义 $y$ 的理想. 可检验
\[
\delta^2\bar{\mathfrak q}\cap\mathcal O_{Y,y}
=\delta^2\mathfrak q.
\]
因此, 将标量扩张到 $\bar k$ 后验证 (1.2.11.1) 即可; 可假设 $a=0$.\footnote{印本此处引用 (4.11.1); 正在证明的同余式是 (1.2.11.1).}
在 $y$ 处于是有
\[
\delta^2\mathfrak q=
(x_0^5,(x_0^3x_i)_{1\leq i},
 (x_0x_ix_j)_{1\leq i,j},
 (x_ix_jx_k)_{1\leq i,j,k}).
\]
由此 $\delta^2\mathfrak q\supset\mathfrak q^5$, 并且
\[
x_0^2f\equiv x_0^4\pmod{\delta^2\mathfrak q},
\qquad
\delta^2\mathfrak q+(x_0^2f)\supset\mathfrak q^4,
\]
\[
x_0f\equiv x_0^3
 \pmod{\delta^2\mathfrak q+\mathfrak q^4},
\qquad
\delta^2\mathfrak q+(x_0f)\supset\mathfrak q^3.
\]
最后,
\[
Q(x_i)\equiv0\pmod{\delta^2\mathfrak q+(f)}
\qquad\underline{\text{在 }Z},
\]
这就是 (1.2.11.1).

\subsubsection*{注 1.2.12.}
设 $Q$ 为不必齐次的 $n+1$ 元非零二次型, 并假设由方程 $Q$ 定义的闭子概形 $Y\subset\mathbb A_k^{n+1}$ 在某点 $y$ 处有普通二次奇点. 则经过一个非齐次线性换元, 可将 $Y$ 化为 1.2.3 或 1.2.4 的形式. 非退化时 $k(y)=k$, 只须将 $y$ 移到原点. 退化时, 只须使 $T^1_{Y/k}$ 的支集落在直线 $x_i=0$ $(1\leq i\leq2m)$ 上(参见 1.2.9).